\section*{Chapitre \ref{ch:b2:blood-circulation} --- Le sang et l'appareil circulatoire}

\begin{solution}{exo:b2:blood-circulation:1}
\omterm{def:b2:blood-circulation:blood}{Plasma} \qty{55}{\%} (eau, \qty{70}{g/L} de protéines, sels, nutriments,
gaz, hormones)~; cellules \qty{45}{\%}. \omterm{def:b2:blood-circulation:blood}{Hématies} $5\times
10^{12}$/L, disques de \qty{7}{\micro m}, transport du dioxygène~;
leucocytes
$7\times 10^{9}$/L, \qtyrange{10}{15}{\micro m}, immunité~; \omterm{def:b2:blood-circulation:blood}{plaquettes}
$3\times 10^{11}$/L, fragments de \qty{2}{\micro m}, hémostase.
\end{solution}

\begin{solution}{exo:b2:blood-circulation:2}
Poisson~: cœur $\to$ branchies $\to$ corps $\to$ cœur, un seul circuit,
le corps recevant le \omterm{def:b2:blood-circulation:blood}{sang} à la faible pression qui subsiste après les
branchies. Mammifère~:
cœur droit $\to$ poumons (\qty{25}{mmHg}) $\to$ cœur gauche $\to$ corps
(\qty{120}{mmHg}) $\to$ cœur droit. La circulation double donne au corps
une haute pression sans y exposer les poumons, et permet de régler les
deux circuits séparément.
\end{solution}

\begin{solution}{exo:b2:blood-circulation:3}
\omterm{def:b2:blood-circulation:vessels}{Artère}~: paroi épaisse à couches élastiques et musculaires --- conduit le
\omterm{def:b2:blood-circulation:blood}{sang} à haute pression et amortit les à-coups. \omterm{def:b2:blood-circulation:vessels}{Artériole}~: paroi surtout
faite de muscle lisse --- fixe la résistance et répartit le débit.
\omterm{def:b2:blood-circulation:vessels}{Capillaire}~: paroi épaisse d'une seule cellule endothéliale --- échanges.
\omterm{def:b2:blood-circulation:vessels}{Veine}~: paroi mince et distensible, pourvue de valvules --- ramène le \omterm{def:b2:blood-circulation:blood}{sang}
à basse pression et en stocke la plus grande partie.
\end{solution}

\begin{solution}{exo:b2:blood-circulation:4}
$70\times 72 = \qty{5}{L/min}$, \qty{7200}{L} par jour, pour \qty{5}{L}
de \omterm{def:b2:blood-circulation:blood}{sang}~: ce volume passe par le cœur \num{1440} fois par jour. Aucun
organe ne pourrait fabriquer ni consommer sept tonnes de \omterm{def:b2:blood-circulation:blood}{sang} par jour~;
le même \omterm{def:b2:blood-circulation:blood}{sang} doit revenir --- il circule.
\end{solution}

\begin{solution}{exo:b2:blood-circulation:5}
$\Delta P = 8\eta LQ/\pi r^{4} = 8\times 3\times 10^{-3}\times
0.4\times 8.3\times 10^{-5}/(\pi\times 2.44\times 10^{-8}) =
\qty{10}{Pa}$, soit \qty{0.08}{mmHg}~: l'aorte ne coûte rien~; la
pression est dépensée dans les \omterm{def:b2:blood-circulation:vessels}{artérioles}.
\end{solution}

\begin{solution}{exo:b2:blood-circulation:6}
$(15/12)^{4} = 2.4$~: résistance multipliée par 2.4, débit ramené à
\qty{41}{\%}. $(15/20)^{4} = 0.32$~: résistance ramenée au tiers, débit
multiplié par 3.2.
\end{solution}

\begin{solution}{exo:b2:blood-circulation:7}
Aorte $83/3 = \qty{28}{cm/s}$~; \omterm{def:b2:blood-circulation:vessels}{capillaires} $83/3000 =
\qty{0.028}{cm/s}$~; transit $0.1/0.028 = \qty{3.6}{s}$. À l'effort, le
débit est multiplié par cinq et la section des \omterm{def:b2:blood-circulation:vessels}{capillaires} ouverts par
trois environ~: le temps de transit tombe à environ \qty{2}{s} --- encore
assez pour que le dioxygène sorte.
\end{solution}

\begin{solution}{exo:b2:blood-circulation:8}
Extrémité artérielle $35 - 25 = \qty{+10}{mmHg}$ (filtration)~;
extrémité veineuse $15
- 25 = \qty{-10}{mmHg}$ (réabsorption). Avec $P_c = 28$ à l'extrémité
veineuse, le bilan est de $+3$~: le liquide filtre sur toute la longueur
et rien n'est réabsorbé --- c'est l'œdème, les chevilles gonflées de
l'insuffisance cardiaque.
\end{solution}

\begin{solution}{exo:b2:blood-circulation:9}
$40/66000 = \qty{0.61}{mol/m^3}$~; $\pi = 0.61\times 8.314\times 310 =
\qty{1570}{Pa} = \qty{11.8}{mmHg}$~; pour moitié moins d'albumine,
\qty{5.9}{mmHg}. Le sodium traverse librement la paroi \omterm{def:b2:blood-circulation:vessels}{capillaire}~: sa
concentration est la même des deux côtés et il n'exerce aucune pression
osmotique à travers elle.
\end{solution}

\begin{solution}{exo:b2:blood-circulation:10}
$RC = \qty{2}{s}$~: $120e^{-0.4} = \qty{80}{mmHg}$. $C = 1$~: $RC = 1$,
$120e^{-0.8} = \qty{54}{mmHg}$ --- et le même volume d'éjection élève
davantage la pression systolique dans une aorte rigide. Une \omterm{thm:b2:blood-circulation:windkessel}{pression
différentielle} élevée signifie un pic plus haut contre lequel le cœur
doit éjecter, donc plus de travail et plus de dioxygène pour le même
débit.
\end{solution}

\begin{solution}{exo:b2:blood-circulation:11}
Repos~: $90/83 = \qty{1.08}{mmHg.s/mL}$~; effort~: $105/417 =
\qty{0.25}{mmHg.s/mL}$, une baisse d'un facteur quatre. Comme
$R \propto r^{-4}$, $r$ a été multiplié par $4^{1/4} = 1.41$~: les
\omterm{def:b2:blood-circulation:vessels}{artérioles} se sont élargies de \qty{41}{\%} en moyenne.
\end{solution}

\begin{solution}{exo:b2:blood-circulation:12}
La continuité fixe la vitesse à chaque niveau par la section totale, et
l'arbre est construit pour que cette section soit mille fois celle de
l'aorte dans les \omterm{def:b2:blood-circulation:vessels}{capillaires} et redevienne petite dans les \omterm{def:b2:blood-circulation:vessels}{veines}~; la
\omterm{thm:b2:blood-circulation:poiseuille}{loi de Poiseuille} place la résistance dans les \omterm{def:b2:blood-circulation:vessels}{artérioles}, où le muscle
peut la modifier, et laisse les gros vaisseaux presque sans perte. Une
\omterm{def:b2:blood-circulation:circuits}{circulation ouverte} n'a pas de \omterm{def:b2:blood-circulation:vessels}{capillaires} pour ralentir le \omterm{def:b2:blood-circulation:blood}{sang} en un
lieu défini, pas de vaisseaux où fixer une résistance, et aucun moyen
de diriger le débit vers un organe --- elle ne peut que brasser.
\end{solution}

\begin{solution}{pb:b2:blood-circulation:1}
\textbf{1.} $70\times 72 = \qty{5.0}{L/min}$~; \qty{7300}{L} par jour.
\textbf{2.} $7300/5 \approx 1450$ fois.
\textbf{3.} Le quart de \qty{57}{g} à 72 battements par minute~:
\qty{14}{g}
$\times 4320 = \qty{62}{kg}$ par heure --- à peu près le poids d'un homme.
\textbf{4.} Aucune nourriture ne pourrait fournir, et aucun tissu
consommer, le poids d'un homme en \omterm{def:b2:blood-circulation:blood}{sang} chaque heure~; la seule issue est
que le même \omterm{def:b2:blood-circulation:blood}{sang} revienne au cœur.
\textbf{5.} Un lien serré autour du bras fait gonfler les \omterm{def:b2:blood-circulation:vessels}{veines} situées
en dessous, non au-dessus~: le \omterm{def:b2:blood-circulation:blood}{sang} des \omterm{def:b2:blood-circulation:vessels}{veines} se dirige donc vers le
cœur~; si l'on pousse le \omterm{def:b2:blood-circulation:blood}{sang} d'une \omterm{def:b2:blood-circulation:vessels}{veine} vers la main, il s'arrête aux
valvules et ne peut pas être refoulé --- les valvules ne permettent
l'écoulement que vers le cœur.
\textbf{6.} Les \omterm{def:b2:blood-circulation:vessels}{capillaires} qui relient les \omterm{def:b2:blood-circulation:vessels}{artères} aux \omterm{def:b2:blood-circulation:vessels}{veines}~;
Malpighi les observa dans le poumon de la grenouille en 1661.
\textbf{7.} Section $\pi\times 1.25^{2} = \qty{4.9}{cm^2}$~; $83/4.9 =
\qty{17}{cm/s}$.
\textbf{8.} $8\times 3\times 10^{-3}\times 0.4\times 8.3\times
10^{-5}/(\pi\times 2.44\times 10^{-8}) = \qty{10}{Pa}$, soit
\qty{0.08}{mmHg}.
\textbf{9.} $R_{1} = 8\times 3\times 10^{-3}\times 10^{-3}/(\pi\times
5.06\times 10^{-20}) = \qty{1.5e14}{Pa.s/m^3}$~; en parallèle,
$1.5\times 10^{14}/3\times 10^{6} = \qty{5.0e7}{Pa.s/m^3}$.
\textbf{10.} $8.3\times 10^{-5}\times 5.0\times 10^{7} = \qty{4200}{Pa}
= \qty{31}{mmHg}$.
\textbf{11.} \omterm{def:b2:blood-circulation:vessels}{Capillaires} ouverts~: $10^{10}$, chacun de section
$\pi(3\times 10^{-6})^{2}
= \qty{2.8e-11}{m^2}$~: \qty{0.28}{m^2} $= \qty{2800}{cm^2}$~; $v =
83/2800 = \qty{0.03}{cm/s}$~; transit $0.1/0.03 = \qty{3.3}{s}$.
\textbf{12.} Résistance $\times(1/0.9)^{4} = 1.52$~: \qty{47}{mmHg}~; le
cœur doit élever la pression artérielle de \qty{16}{mmHg}, faute de quoi
le débit chute d'un tiers.
\textbf{13.} $+10$, $-10$ et $0$ mmHg.
\textbf{14.} Filtration $10\times 2 = \qty{20}{L}$ par jour~;
réabsorption $8.5\times 2 = \qty{17}{L}$.
\textbf{15.} \qty{3}{L} de lymphe par jour.
\textbf{16.} $\pi_c = \qty{12.5}{mmHg}$~: bilan $+22.5$ à l'extrémité
artérielle et $+2.5$ à l'extrémité veineuse --- filtration partout, pas
de réabsorption~: œdème.
\textbf{17.} Bilans $+10$ et $+3$, moyenne $+6.5$~: filtration sur toute
la longueur du \omterm{def:b2:blood-circulation:vessels}{capillaire}, quelque \qty{26}{L} par jour~; les
lymphatiques ne peuvent pas l'évacuer et le liquide s'accumule dans les
tissus, en commençant par les pieds.
\textbf{18.} $4\times 10^{10}\times 2\pi\times 3\times 10^{-6}\times
10^{-3} = \qty{750}{m^2}$~; $t = x^{2}/2D = (2\times 10^{-5})^{2}/
(2\times 10^{-9}) = \qty{0.2}{s}$.
\textbf{19.} $(93 - 3)/83 = \qty{1.08}{mmHg.s/mL}$.
\textbf{20.} $RC = \qty{2.2}{s}$~; $120e^{-0.8/2.2} = 120\times 0.69 =
\qty{83}{mmHg}$.
\textbf{21.} $RC = \qty{1.1}{s}$~; $120e^{-0.73} = \qty{58}{mmHg}$~: la
\omterm{thm:b2:blood-circulation:windkessel}{pression différentielle} passe de 37 à \qty{62}{mmHg}.
\textbf{22.} $120e^{-0.3/2.2} = \qty{105}{mmHg}$~: à fréquence élevée, la
pression diminue à peine entre deux battements.
\textbf{23.} $C\Delta P = 2\times 20 = \qty{40}{mL}$ stockés~; les
\qty{30}{mL} restants s'écoulent par les \omterm{def:b2:blood-circulation:vessels}{artérioles} pendant la systole
elle-même.
\textbf{24.} Le ventricule qui se contracte écrase ses propres vaisseaux
en systole~: le muscle cardiaque est donc irrigué en diastole, grâce à
la pression que maintient la rétraction de l'aorte~; une aorte rigide
qui laisse s'effondrer la pression diastolique prive le cœur de \omterm{def:b2:blood-circulation:blood}{sang}
entre deux battements.
\textbf{25.} Débit \qty{5}{L/min}~; chute artériolaire \qty{31}{mmHg}~;
filtrat net \qty{3}{L} par jour vers la lymphe~; pression diastolique
\qty{83}{mmHg} pour $C = 2$ et \qty{58}{mmHg} pour $C = 1$.
\end{solution}
